\section{Summary}
\label{sec:precise-alignment-summary}

This chapter presented a GICP-based point cloud registration pipeline for centimetre-scale pose refinement against a 3D Gaussian Splatting reference model, evaluated with two sensing modalities across a $3 \times 3$ factorial design of sensor-to-vine distance and approach angle. The LiDAR pipeline achieved 100\% condition-level convergence (311/315 hypotheses accepted) with a median inlier RMSE of 33.25\,mm and mean latency of 46\,s, while the stereo pipeline converged in seven of nine conditions with a mean inlier RMSE of 15.51\,mm and mean latency of 14.18\,s. Neither modality dominated: LiDAR provided superior convergence reliability attributable to its active sensing physics, whereas stereo achieved approximately half the RMSE when converging. The discussion highlighted that the multi-hypothesis perturbation strategy provided redundancy rather than diversity, that moderate operating distance (60--80\,cm) produced the best registration quality for both modalities, and that the ICP corrections fell within the SLAM system's expected error bounds, validating the two-stage localisation architecture.